% Auto-generated by scripts/aggregate_section_5_6.py
% Theorem-receipts summary table — one row per formal claim with verdict + receipt anchor.
\begin{table}[!htbp]
\centering\scriptsize
\setlength{\tabcolsep}{4pt}
\renewcommand{\arraystretch}{0.95}
\begin{tabular}{@{}p{0.4cm}p{3.0cm}p{1.1cm}p{1.7cm}p{5.4cm}p{2.4cm}@{}}
\toprule
\textbf{ID} & \textbf{Label} & \textbf{Type} & \textbf{Verdict} & \textbf{Receipt summary} & \textbf{Anchor} \\
\midrule
T1 & \texttt{thm:\allowbreak purity} & Theorem & PASS (pooled cal-fold) & Pooled cal-fold $\pi_t \in [0.732, 0.746]$ clears the $0.64$ floor by $9$ to $11$ percentage points at every reported cycle (the registered pooled training-panel axis). Per-bench eval-fold FEVER dips below floor at C1--C4 (band $[0.614, 0.634]$); disclosed as within-scope-permitted diagnostic, outside the pooled-axis bound. & \Cref{tab:pi-store-trajectory} \\
T2 & \texttt{thm:\allowbreak bayes-\allowbreak purity} & Remark & identity & Mathematical scaffolding (Bayes'-rule identity); no within-horizon empirical receipt. & -- \\
T3 & \texttt{thm:\allowbreak monotone} & Theorem & consistent within horizon, sharp test deferred & Sign alignment holds on three of four cal-fold transitions; honest reverse-sign disclosure at $C3 \to C5$. & \Cref{tab:pi-store-trajectory} \\
T4 & \texttt{thm:\allowbreak monotone-\allowbreak purification} & Remark & identity & Mathematical scaffolding (derivative of T2); no within-horizon empirical receipt. & -- \\
T5 & \texttt{thm:\allowbreak convergence} & Theorem & BOUNDED-BAND & Pooled $\pi_t$ confined to a band of width $\le 0.031$ across the six-cycle horizon; the rate-scaling form $r = \mathcal{O}\!\big(\sigma \cdot \Phi(d/\sqrt{2})\big)$ is a Gaussian-model heuristic asserted in the proof sketch rather than measured at the present horizon, and the measured geometric-rate test is deferred to a $\ge 15$-committed-cycle horizon. & \Cref{tab:pi-store-trajectory} \\
T6 & \texttt{thm:\allowbreak bayes-\allowbreak convergence} & Remark & identity & Mathematical scaffolding (monotone-map fixed-point theorem); no within-horizon empirical receipt. & -- \\
T7 & \texttt{thm:\allowbreak asymptotic-\allowbreak elim} & Remark & identity & Mathematical scaffolding (law of total probability); no within-horizon empirical receipt. & -- \\
C1 & \texttt{cor:\allowbreak tier1-\allowbreak floor} & Corollary & CAPABILITY & CHM$_1 = 0$ on each of the seven realised Tier-1 firings; pool-side ceiling $1 - \pi_{\text{store}} \approx 0.36$ holds with substantial slack. & \Cref{tab:tier-chm-trajectory} \\
C3 & \texttt{cor:\allowbreak corpus-\allowbreak floor} & Remark & deferred & Per-sample top-$k$ passage IDs not persisted; proxy receipt deferred to future-work re-retrieval pass. & -- \\
C4 & \texttt{cor:\allowbreak self-\allowbreak correction} & Corollary & PARTIAL & Prune-rate decline $13 \to 9 \to 7 \to 4$ percent across committed cycles supports directional claim; survivor-side histogram deferred. & \Cref{tab:self-correction-check} \\
C5 & \texttt{cor:\allowbreak stochastic-\allowbreak equilibrium} & Corollary & PASS (i,iii) / heuristic (iv) & Sub-claims (i) and (iii) pass directly: $\pi_{\mathrm{commit}} = 4/5 = 0.80$ strictly in $(0,1)$; memory grows monotonically from $431$ to $12\,175$ entries across the C4 abort. Sub-claim (ii) inherits the bounded-band reading of T5. Sub-claim (iv), the equilibrium-location prediction, is labelled heuristic at the five-cycle horizon and deferred to a $\ge 15$-committed-cycle test. & \Cref{tab:stochastic-equilibrium} \\
\bottomrule
\end{tabular}
\caption{Theorem and corollary receipts summary against the realised six-cycle trajectory. Each row pairs one formal claim from \Cref{sec:theory} with a verdict and the table anchor where its empirical receipt is reported. PASS verdicts are claims whose receipt is satisfied directly by the realised trajectory. PARTIAL verdicts are claims whose directional prediction is supported by the realised data but whose sharper form requires either a longer horizon or an instrument not present in the artefact set. BOUNDED-BAND, PROJECTED, and CAPABILITY verdicts label claims that the chapter-four restructuring carries with explicit scope qualifiers as discussed in \Cref{sec:theory}. The C5 row reports a per-sub-claim verdict because the corollary has four sub-claims of different horizons: sub-claims (i) and (iii) pass within the realised five-cycle trajectory, sub-claim (ii) inherits T5's bounded-band reading, and sub-claim (iv) is the equilibrium-location prediction the proof labels heuristic, whose rigorous test requires the longer-horizon item registered in \Cref{ch:conclusion}. Mathematical scaffolding identities (Remarks) do not carry an empirical receipt by design and are listed for completeness.}
\label{tab:theorem-receipts-summary}
\end{table}
